\section{Literature Review}

\subsection{AI-Enabled Demand Planning and Forecasting Progress}

Recent supply-chain research shows rapid progress in AI-enabled demand planning. Walter et al. (2025) identify AI tools, supply-chain applications, and digital transformation as major knowledge clusters; Culot et al. (2024) survey empirical AI adoption and emphasize connecting AI advancement to operational outcomes; Vlachos and Reddy (2025) identify decision-support, execution optimization, and control monitoring as high-value areas. Forecasting benchmarks show the same trend: the M5 competition demonstrates modern retail forecasting at scale across hierarchies and intermittency (Makridakis et al. 2022), and recent time-series foundation models suggest forecasting pipelines are moving toward broader zero-shot and cross-domain capability (Das et al. 2024; Ansari et al. 2024). These advances make deployment evaluation more urgent: a more accurate forecast is valuable only if operations can execute the resulting plan.

\subsection{Forecast-to-Decision and Predict-Then-Optimize Research}

Forecast-to-decision research directly addresses why prediction accuracy and downstream decision quality can diverge. Elmachtoub and Grigas (2022) show predict-then-optimize workflows improve by training toward decision loss rather than prediction loss alone; Mandi et al. (2024) survey decision-focused learning, integrating machine learning and constrained optimization to improve decision quality under uncertainty; and Jia et al. (2025) propose scenario predict-then-optimize for online inventory routing, showing routing and replenishment decisions need more than point-forecast accuracy. This paper shares that motivation but a different goal: rather than new forecasting or end-to-end learning algorithms, it proposes an interpretable evaluation layer that lets operations teams compare existing candidate forecasts by operational consequence, not accuracy metrics alone.

\subsection{Operations Management and Planning-System Architecture}

Advanced planning and ERP systems architecturally separate forecasting, planning, and execution layers (Stadtler et al. 2015), and model deployment requires validation, approval, and integration cycles that take weeks or months---so when forecasting models change faster than implementation can absorb, operations teams face a governance problem, not just a technical one. Inventory theory provides the operational foundation: Silver et al. (1998) and Axsater (2015) frame inventory planning as a cost-and-service-level problem rather than a pure prediction problem, Hyndman and Koehler (2006) show forecast-error measures have known limitations, and forecast combination can reduce model-specific swings (Bates and Granger 1969). Forrester (1961) shows that industrial systems amplify instability when material and information flows respond with delays---the bullwhip effect---and the same logic applies here: coordination mechanisms (supplier partnerships, EDI, collaborative planning) reduce the planning-execution gap but do not eliminate it, because governance processes are slow to change compared to forecasting innovation. Rolling-horizon decision theory adds the final piece: Sethi and Sorger (1991) show a locally attractive action can be undesirable once future transitions are considered, which motivates this paper's greedy-to-DP comparison as a measure of planning myopia.

\subsection{Research Gap and This Paper's Position}

Existing research improves forecasts, combines forecasts, or optimizes downstream decisions in isolation. Fewer studies provide a practical, interpretable decision layer that compares forecasting-driven strategies by their operational consequences---inventory exposure, plan stability, execution violations, governance burden---directly, at the forecast-to-plan boundary where execution actually happens. This paper addresses that gap; the goal is not to replace forecast accuracy as a criterion but to position it as one signal among several operational objectives, evaluated where planning decisions are executed rather than where they are generated.

